% Options for packages loaded elsewhere
\PassOptionsToPackage{unicode}{hyperref}
\PassOptionsToPackage{hyphens}{url}
\PassOptionsToPackage{dvipsnames,svgnames,x11names}{xcolor}
%
\documentclass[
  11pt,
]{article}
\usepackage{amsmath,amssymb}
\usepackage{iftex}
\ifPDFTeX
  \usepackage[T1]{fontenc}
  \usepackage[utf8]{inputenc}
  \usepackage{textcomp} % provide euro and other symbols
\else % if luatex or xetex
  \usepackage{unicode-math} % this also loads fontspec
  \defaultfontfeatures{Scale=MatchLowercase}
  \defaultfontfeatures[\rmfamily]{Ligatures=TeX,Scale=1}
\fi
\usepackage{lmodern}
\ifPDFTeX\else
  % xetex/luatex font selection
\fi
% Use upquote if available, for straight quotes in verbatim environments
\IfFileExists{upquote.sty}{\usepackage{upquote}}{}
\IfFileExists{microtype.sty}{% use microtype if available
  \usepackage[]{microtype}
  \UseMicrotypeSet[protrusion]{basicmath} % disable protrusion for tt fonts
}{}
\makeatletter
\@ifundefined{KOMAClassName}{% if non-KOMA class
  \IfFileExists{parskip.sty}{%
    \usepackage{parskip}
  }{% else
    \setlength{\parindent}{0pt}
    \setlength{\parskip}{6pt plus 2pt minus 1pt}}
}{% if KOMA class
  \KOMAoptions{parskip=half}}
\makeatother
\usepackage{xcolor}
\usepackage[left=1.25in,right=1.25in,top=1in,bottom=1in]{geometry}
\usepackage{longtable,booktabs,array}
\usepackage{calc} % for calculating minipage widths
% Correct order of tables after \paragraph or \subparagraph
\usepackage{etoolbox}
\makeatletter
\patchcmd\longtable{\par}{\if@noskipsec\mbox{}\fi\par}{}{}
\makeatother
% Allow footnotes in longtable head/foot
\IfFileExists{footnotehyper.sty}{\usepackage{footnotehyper}}{\usepackage{footnote}}
\makesavenoteenv{longtable}
\setlength{\emergencystretch}{3em} % prevent overfull lines
\providecommand{\tightlist}{%
  \setlength{\itemsep}{0pt}\setlength{\parskip}{0pt}}
\setcounter{secnumdepth}{-\maxdimen} % remove section numbering
\ifLuaTeX
  \usepackage{selnolig}  % disable illegal ligatures
\fi
\IfFileExists{bookmark.sty}{\usepackage{bookmark}}{\usepackage{hyperref}}
\IfFileExists{xurl.sty}{\usepackage{xurl}}{} % add URL line breaks if available
\urlstyle{same}
\hypersetup{
  pdftitle={Follow-up note: empirical triangulation of (r\_e/r\_0)},
  pdfauthor={Trey Morris with Claude Opus 4.7},
  pdfsubject={Empirical triangulation of r\_e/r\_0 across six precision
observables (follow-up to interim report)},
  colorlinks=true,
  linkcolor={blue},
  filecolor={Maroon},
  citecolor={Blue},
  urlcolor={blue},
  pdfcreator={LaTeX via pandoc}}

\title{Follow-up note: empirical triangulation of (r\_e/r\_0)}
\author{Trey Morris with Claude Opus 4.7}
\date{2026-06-01}

\begin{document}
\maketitle

\hypertarget{follow-up-note-empirical-triangulation-of-r_er_0-for-tepper-gill}{%
\section{\texorpdfstring{Follow-up note: empirical triangulation of
\(r_e/r_0\) --- for Tepper
Gill}{Follow-up note: empirical triangulation of r\_e/r\_0 --- for Tepper Gill}}\label{follow-up-note-empirical-triangulation-of-r_er_0-for-tepper-gill}}

\textbf{Date:} 2026-05-26. \textbf{From:} Trey Morris (with Claude Opus
4.7). \textbf{Re:} Author guidance recorded after the 2026-05-25 phone
call following the interim report\footnote{Interim report, pull request
  \#59 (\texttt{github.com/temoTxt/PyPhysics/pull/59}). §5 of that
  report sketches the branches (b) and (c) discussed here.}; empirical
triangulation closes the joint-fit thread\footnote{Issue \#61
  (\texttt{github.com/temoTxt/PyPhysics/issues/61}). Empirical
  triangulation thread.}; the first-principles rederivation continues
separately\footnote{Issue \#54
  (\texttt{github.com/temoTxt/PyPhysics/issues/54}). First-principles
  rederivation thread.}.

\begin{center}\rule{0.5\linewidth}{0.5pt}\end{center}

Tepper --- following your guidance that the branches (b) and (c) in §5
of the interim report\footnote{Interim report, pull request \#59
  (\texttt{github.com/temoTxt/PyPhysics/pull/59}). §5 of that report
  sketches the branches (b) and (c) discussed here.} are bracketing
guides from a uni-observable search rather than theoretical predictions,
we have performed the joint fit across all six precision observables
that you suggested. We summarise here what we found. The full notebook
and per-observable diagnostics are recorded separately\footnote{\texttt{Roadmapping/Mathematica\_Notebooks/Quantum\_Mechanics/r\_e\_triangulation.wl}.
  Joint-fit notebook with per-observable diagnostics.}.

\hypertarget{triangulated-value}{%
\subsection{1. Triangulated value}\label{triangulated-value}}

We obtain, from a precision-weighted \(\chi^2\) joint fit,

\[
r_e/r_0 \;=\; 0.499\,420\,509\,912\,831\,7 \;\pm\; 2.5\times 10^{-13},
\qquad (1)
\]

where the uncertainty is the one-sigma value taken from the second
derivative of \(\chi^2\) at the optimum.

We observe that (1) differs from the published
\(r_e/r_0 = 0.499\,857\,150\,068\,631\) in the fourth decimal place. It
is, however, indistinguishable from the back-fit-to-\(g_s\) value
(branch (c) of the interim report) to sixteen significant figures. The
triangulation thus makes quantitative what the Bethe--Salpeter campaign
already documented structurally\footnote{\texttt{Roadmapping/Quantum\_Mechanics/Bethe\_Salpeter/10\_CrossComparison.md},
  §2 (the \(r_e\) back-fit self-consistency across six \(g_s\)-dependent
  observables).}. Indeed, every \(g_s\)-dependent observable in the
framework satisfies

\[
\text{prediction} \;=\; (g_s / -2)^n \times \text{textbook}, \qquad n \in \{1, 2\},
\qquad (2)
\]

so a single back-fit applied to six observables yields a single \(r_e\)
value.

\hypertarget{per-observable-residuals-at-the-triangulated-optimum}{%
\subsection{2. Per-observable residuals at the triangulated
optimum}\label{per-observable-residuals-at-the-triangulated-optimum}}

The fit residuals are collected in Table 1.

\begin{longtable}[]{@{}
  >{\raggedright\arraybackslash}p{(\columnwidth - 8\tabcolsep) * \real{0.2000}}
  >{\raggedright\arraybackslash}p{(\columnwidth - 8\tabcolsep) * \real{0.2000}}
  >{\raggedright\arraybackslash}p{(\columnwidth - 8\tabcolsep) * \real{0.2000}}
  >{\raggedright\arraybackslash}p{(\columnwidth - 8\tabcolsep) * \real{0.2000}}
  >{\raggedright\arraybackslash}p{(\columnwidth - 8\tabcolsep) * \real{0.2000}}@{}}
\toprule\noalign{}
\begin{minipage}[b]{\linewidth}\raggedright
Observable
\end{minipage} & \begin{minipage}[b]{\linewidth}\raggedright
Prediction
\end{minipage} & \begin{minipage}[b]{\linewidth}\raggedright
Measurement
\end{minipage} & \begin{minipage}[b]{\linewidth}\raggedright
Residual
\end{minipage} & \begin{minipage}[b]{\linewidth}\raggedright
Source of residual
\end{minipage} \\
\midrule\noalign{}
\endhead
\bottomrule\noalign{}
\endlastfoot
Electron \(g_s\) & \(-2.00231930\ldots\) & \(-2.00231930\ldots\) & \(0\)
& matches by construction \\
H \(2P_{3/2}-2P_{1/2}\) & \(10{,}962\) MHz & \(10{,}969.13\) MHz &
\(-7.13\) MHz & Bethe-estimate floor\footnote{\texttt{Roadmapping/Quantum\_Mechanics/Bethe\_Salpeter/14\_Lamb\_Shift\_2P.md},
  §14.2 (Bethe-estimate precision floor).} \\
H 1S hyperfine & \(1{,}420.04\) MHz & \(1{,}420.4058\) MHz & \(-0.366\)
MHz & sub-leading QED\footnote{\texttt{Roadmapping/Quantum\_Mechanics/Bethe\_Salpeter/22\_Hydrogen\_Hyperfine.md},
  §22.1 (sub-leading QED gap).} \\
He \({}^3P_0 - {}^3P_1\) & \(29{,}616.95\) MHz & \(29{,}616.952\) MHz &
\(-0.002\) MHz & kHz floor\footnote{\texttt{Roadmapping/Quantum\_Mechanics/Bethe\_Salpeter/07\_Helium\_Fine\_Structure.md},
  §72 (kHz precision floor).} \\
Positronium ortho-para & \(203{,}389\) MHz & \(203{,}389\) MHz & \(0\) &
matches by construction \\
Muonium hyperfine & \(4{,}463.40\) MHz & \(4{,}463.3028\) MHz &
\(+0.097\) MHz & sub-leading QED\footnote{\texttt{Roadmapping/Quantum\_Mechanics/Bethe\_Salpeter/08\_Muonium\_Hyperfine.md},
  §80 (sub-leading QED gap).} \\
\end{longtable}

\emph{Table 1.} Per-observable residuals at the triangulated optimum
(1).

Every non-zero residual in Table 1 is a documented
framework-precision-floor effect. It is not reducible by re-tuning
\(r_e/r_0\). No observable is in tension with (1) once the framework's
known precision scope is respected.

\hypertarget{honest-scoping-note-on-the-objective-function}{%
\subsection{3. Honest scoping note on the objective
function}\label{honest-scoping-note-on-the-objective-function}}

We record one substantive-AI decision exposed by the joint fit. The
literal ``fit to all six measurements'' reading treats each
measurement's stated uncertainty as the \(\sigma_i\) entering
\(\chi^2\). Under that reading the optimum is

\[
r_e/r_0 \;=\; 0.499\,406\,1\ldots,
\qquad (3)
\]

differing from (1) in the fifth decimal place. Equation (3) pulls
\(g_e\) off the measured value by approximately

\[
\Delta g_e \;\approx\; 5.8 \times 10^{-5},
\qquad (4)
\]

in an attempt to compensate the hyperfine and helium fine-structure
residuals. This is an artifact of treating framework-precision-floor
residuals as if they were measurement uncertainty.

The honest formulation adds a framework-floor noise term to each
\(\sigma_i\),

\[
\sigma_i^{2} \;=\; \sigma_{\text{meas},\,i}^{2} \;+\; \sigma_{\text{floor},\,i}^{2}.
\qquad (5)
\]

The recognition behind (5) is that the leading \((g_s/-2)^n\) model
cannot fit Bethe-estimate-gap residuals by any choice of \(r_e\) alone.
Under the weighting prescribed by (5) the optimum is the value (1),
matching branch (c) of the interim report to sixteen significant
figures. Both passes are documented in the notebook\footnote{\texttt{Roadmapping/Mathematica\_Notebooks/Quantum\_Mechanics/r\_e\_triangulation.wl}.
  Joint-fit notebook with per-observable diagnostics.}. We send you the
honest answer here, with the alternative (3) recorded for transparency.

\hypertarget{where-this-leaves-the-open-question}{%
\subsection{4. Where this leaves the open
question}\label{where-this-leaves-the-open-question}}

A natural way to read the relationship between the published \(r_e\) and
the triangulated value follows. The published
\(r_e/r_0 = 0.499\,857\ldots\) is an \emph{initial-value} result from
the working-notebook numerical search against \(g_s\) alone. The joint
fit across six observables is a \emph{refinement calculation} on top of
that initial value. The triangulated value (1) supersedes the initial
value in the same way that any iteratively-refined numerical result
supersedes an earlier estimate. It does not invalidate the derivation
that produced the initial value. It refines the cutoff identification
using six constraints rather than one.

A further refinement is still open: a first-principles derivation of
\(r_e\) from the dual-Dirac renormalisation prescription, in the sense
of DRQM I §III.D. This is tracked separately\footnote{Issue \#54
  (\texttt{github.com/temoTxt/PyPhysics/issues/54}). First-principles
  rederivation thread.}. The candidate starting points are sketched
briefly below, and expanded into one-to-three-paragraph explanations
each in a companion note\footnote{\texttt{Roadmapping/Author\_Reports/2026-05\_re\_derivation\_candidates\_for\_gill.md}.
  Companion note expanding the three candidate starting points.}.

By ``first-principles'' we mean a derivation in which \(r_e\) emerges
from the dual-Dirac equation's internal structure, rather than being
fixed by inverse-solving against the measured \(g_e\). The output would
be a closed-form expression for \(r_e/r_0\) in terms of the framework's
fundamental parameters --- functions of \(\alpha\), \(r_0\), and the
structural constants of the dual representation. Such an expression
would evaluate numerically to approximately \(0.499\,420\,509\ldots\) if
the empirical refinement holds, or to a different value if the
derivation exposes a structure the triangulation has not captured.

We have considered three plausible starting points, and would defer to
your judgment on which (if any) is natural for this framework.

\begin{enumerate}
\def\labelenumi{\arabic{enumi}.}
\tightlist
\item
  The proper-time self-energy integral in the dual-Dirac formalism, with
  \(r_e\) emerging as the regulator scale at which mass renormalisation
  closes consistently.
\item
  A variational determination, with \(r_e\) fixed by demanding that the
  renormalised dual-Dirac equation reproduce the physical electron mass
  at the cutoff. This is analogous to the way the standard-QED running
  coupling is fixed by a renormalisation condition.
\item
  A structural constant of the dual representation --- something that
  the \(b\)-factor projection structure or the second-order Dirac
  decomposition singles out --- making \(r_e/r_0\) a derivable
  framework-internal ratio rather than a renormalisation choice.
\end{enumerate}

Per your 2026-05-25 author guidance, a fourth candidate we had initially
considered --- the retrieval of an unpublished working-notebook
derivation step --- is empty. The original cutoff value was a numerical
search against \(g_s\) alone, with no derivation-step working draft to
retrieve.

We have not pursued any of the three. The natural starting point depends
on the framework's internal logic better than we can guess at from
outside. The thread is non-urgent. The triangulated value (1) already
serves as the campaign's current best refinement, and a first-principles
derivation can sit indefinitely as a would-be-nice-to-have until you
point us at a starting structure, or decide the thread is not where the
project's time is best spent.

Whichever route turns out to be natural, the outcome could agree with
the triangulated value (confirming this refinement), supersede it (a
third refinement), or expose a derivation-level structure that reframes
the cutoff entirely. Each of these outcomes is consistent with treating
each calculation as an iterative refinement rather than as a final
authoritative answer.

\hypertarget{references}{%
\subsection{References}\label{references}}

\end{document}
